\chapter*{Annexe C -- Configuration Prometheus}
\addcontentsline{toc}{chapter}{Annexe C -- Configuration Prometheus}
\label{annexe-prometheus}

L'extrait ci-dessous illustre une configuration standard de \textbf{Prometheus} pour la découverte et la supervision des pods d'un cluster Kubernetes, telle qu'évoquée au \cref{chap-sprint4}. \emph{À remplacer par la configuration exacte déployée sur le cluster avant soutenance.}

\begin{lstlisting}[language=bash,caption={Extrait de configuration \texttt{prometheus.yml}}]
global:
  scrape_interval: 15s
  evaluation_interval: 15s

scrape_configs:
  - job_name: 'kubernetes-pods'
    kubernetes_sd_configs:
      - role: pod
    relabel_configs:
      - source_labels: [__meta_kubernetes_pod_annotation_prometheus_io_scrape]
        action: keep
        regex: true
      - source_labels: [__meta_kubernetes_namespace]
        target_label: namespace
      - source_labels: [__meta_kubernetes_pod_name]
        target_label: pod

  - job_name: 'monitoring-backend'
    static_configs:
      - targets: ['backend.monitoring.svc.cluster.local:3000']
\end{lstlisting}

Les métriques ainsi collectées par Prometheus sont ensuite exploitées par \textbf{Grafana} pour la construction de tableaux de bord de supervision de l'infrastructure (état des pods, consommation CPU/mémoire par conteneur, disponibilité des services).
